\section*{Hoofdstuk \ref{ch:g12:destruction-of-mountains} --- De afbraak van gebergten}

\begin{solution}{exo:g12:destruction-of-mountains:1}
(a) de \omterm{def:g12:destruction-of-mountains:denudation}{denudatiesnelheid}; (b) de \omterm{def:g12:destruction-of-mountains:load}{bodemvracht}; (c) de \omterm{def:g12:destruction-of-mountains:load}{opgeloste vracht};
(d) \omterm{def:g12:destruction-of-mountains:diagenesis}{diagenese}; (e) \omterm{def:g12:destruction-of-mountains:rebound}{isostatisch herstel}.
\end{solution}

\begin{solution}{exo:g12:destruction-of-mountains:2}
(1) Een conglomeraat. (2) Een zandsteen. (3) Een schalie. (4) Een breccie.
\end{solution}

\begin{solution}{exo:g12:destruction-of-mountains:3}
\qty{1}{m/Myr} is \qty{1000}{mm} in \num{1000000} jaar, \qty{0.001}{mm} per
jaar. (1) \qty{0.054}{mm} per jaar. (2) \qty{0.218}{mm} per jaar.
(3) $\qty{10}{mm}$ per jaar is \qty{10}{m} in 1000 jaar, \qty{10000}{m/Myr}.
\end{solution}

\begin{solution}{exo:g12:destruction-of-mountains:4}
De vaste vracht is \qty{85}{\percent} van het totaal, dus het totaal is $60/0.85
\approx 70.6$ miljoen ton, en de \omterm{def:g12:destruction-of-mountains:load}{opgeloste vracht} ongeveer $70.6 - 60 \approx
10.6$ miljoen ton per jaar.
\end{solution}

\begin{solution}{exo:g12:destruction-of-mountains:5}
Verwering van het graniet, transport door een rivier, afzetting in een bekken,
begraving onder jongere lagen, cementatie.
\end{solution}

\begin{solution}{exo:g12:destruction-of-mountains:6}
Bij \qty{30}{cm/s} wordt de zandkorrel opgepakt (boven \qty{20}{cm/s}); de
siltkorrel blijft zweven. Bij \qty{5}{cm/s} bezinkt de zandkorrel
(onder \qty{10}{cm/s}), maar de siltkorrel blijft nog altijd zweven, omdat
\qty{0.1}{cm/s} genoeg is om hem zwevend te houden.
\end{solution}

\begin{solution}{exo:g12:destruction-of-mountains:7}
Totale vracht $4/0.85 \approx 4.71$ miljoen ton, $\num{4.71e9}/2830 \approx
\qty{1.66e6}{m^3}$ gesteente per jaar. Over $\qty{5e10}{m^2}$: $\num{3.3e-5}$ m per
jaar, ongeveer \qty{33}{m/Myr}, onder de mediaan van \qty{54}{m/Myr}.
\end{solution}

\begin{solution}{exo:g12:destruction-of-mountains:8}
(1) $1 - 2700/3300 \approx 0.182$. (2) Ongeveer \qty{182}{m}. (3) Iets meer
dan de \qty{163}{m} van het hoofdstuk: hoe lichter de korst vergeleken met de
mantel, hoe minder de kolom stijgt.
\end{solution}

\begin{solution}{exo:g12:destruction-of-mountains:9}
(1) \qty{8}{km} gesteente. (2) $8000 \times 0.163 \approx \qty{1300}{m}$.
(3) $8000/200 = 40$ miljoen jaar.
\end{solution}

\begin{solution}{exo:g12:destruction-of-mountains:10}
(1) $794/43 \approx 18$ miljoen jaar. (2) Het oppervlak zakt maar met
$0.163$ keer de weggeërodeerde dikte, dus het gesteente dat weg moet, is $794/0.163
\approx \qty{4900}{m}$, en dat duurt ongeveer $4900/43 \approx 110$ miljoen jaar.
(3) De continenten worden niet alleen afgesleten: gebergtevorming en het
magmatisme van bogen heffen ze steeds weer op en voegen nieuwe korst toe.
\end{solution}

\begin{solution}{exo:g12:destruction-of-mountains:11}
(1) $\num{1e14}/2830 \approx \qty{3.5e10}{m^3}$, dat is ongeveer
\qty{35}{km^3}. (2) $100/15.5 \approx 6.5$ jaar.
\end{solution}

\begin{solution}{exo:g12:destruction-of-mountains:12}
$E = kh$ en $h = h_0 e^{-fkt}$ geven $E(t) = kh_0 e^{-fkt}$. De dikte die
van $t = 0$ tot $t \to \infty$ wordt weggehaald, is $\int_0^{\infty} kh_0
e^{-fkt}\,\mathrm{d}t = kh_0/(fk) = h_0/f$. Om het oppervlak $h_0$ te verlagen,
moet erosie $h_0/f \approx 6.1\,h_0$ weghalen, zoals de propositie over het herstel
zegt: elke weggeërodeerde meter verlaagt het oppervlak maar met $f$ meter.
\end{solution}

\begin{solution}{exo:g12:destruction-of-mountains:13}
(1) $10 \times 1000 = \qty{10000}{mm}$, \qty{10}{m}. (2) Een ijslast met
dikte $t$ wegnemen is als een massa $\rho_{\text{ijs}}\,t$ per oppervlakte-eenheid
wegnemen; de kolom stijgt tot ze zoveel minder mantel verdringt:
$2000 \times 900/3380 \approx \qty{530}{m}$. (3) Het land stijgt steeds minder
naarmate het zijn nieuwe evenwicht nadert: wat het stijgen aandrijft, is het deel dat
nog niet is hersteld.
\end{solution}

\begin{solution}{exo:g12:destruction-of-mountains:14}
(1) Vóór 1900: ongeveer \qty{56}{m/Myr} (zoals in het voorbeeld). In
1987--2006: $145/0.85 \approx 171$ miljoen ton, $\num{1.71e11}/2830
\approx \qty{6.0e7}{m^3}$, over $\qty{2.98e12}{m^2}$: ongeveer
\qty{20}{m/Myr}. (2) De eerste: de latere vracht wordt tegengehouden door dammen
en andere werken. (3) Nee: de afname is $400 - 145 = 255$ miljoen ton per
jaar, waarvan de dammen 100 tot 150 verklaren, ongeveer de helft.
\end{solution}

\begin{solution}{exo:g12:destruction-of-mountains:15}
(1) Minstens \qty{6}{km}. (2) $310 - 252 = 58$ miljoen jaar hoogstens, dus
minstens $6000/58 \approx \qty{100}{m/Myr}$. (3) Het ligt tussen de mediaan
(\qty{54}{m/Myr}) en het gemiddelde (\qty{218}{m/Myr}) van de \omterm{def:g10:drainage-basins:drainage-basin}{stroomgebieden}: een gewone
snelheid voor een gebergtegordel.
\end{solution}

\begin{solution}{pb:g12:destruction-of-mountains:1}
\textbf{1.} $E = 0.209 \times 2000 \approx \qty{418}{m/Myr}$.

\textbf{2.} \qty{0.418}{mm} per jaar.

\textbf{3.} Sneller dan beide: het is een steile, jonge keten, even snel als de
seismisch actieve \omterm{def:g10:drainage-basins:drainage-basin}{stroomgebieden} en meer.

\textbf{4.} Zijn gemiddelde helling.

\textbf{5.} Een laag van $\qty{0.418e-3}{m}$ over $\qty{1e6}{m^2}$ is
\qty{418}{m^3}, $418 \times 2830 \approx \qty{1.18e6}{kg}$: ongeveer 1200
ton.

\textbf{6.} $2000/418 \approx 4.8$ miljoen jaar.

\textbf{7.} $h(t) = 2000\,e^{-0.209\,t}$ ($t$ in miljoen jaar).

\textbf{8.} $1/0.209 \approx 4.8$ miljoen jaar.

\textbf{9.} $\ln 2/0.209 \approx 3.3$ miljoen jaar.

\textbf{10.} $e^{-0.209 \times 11} = e^{-2.30} \approx 0.100$.

\textbf{11.} $1 - 2830/3380 \approx 0.163$.

\textbf{12.} Het oppervlak zakt $100 \times 0.163 \approx \qty{16}{m}$;
de \omterm{def:g12:continental-crust:root}{wortel} stijgt ongeveer \qty{84}{m}.

\textbf{13.} $h' = -0.163 \times 0.209\,h \approx -0.0341\,h$, dus $h(t) =
2000\,e^{-0.0341\,t}$.

\textbf{14.} $(2000 - 200)/0.163 \approx \qty{11000}{m}$ gesteente.

\textbf{15.} Gesteenten die ongeveer \qty{11}{km} diep lagen: metamorfe gesteenten
zoals schisten en gneisen, en granieten.

\textbf{16.} $1/0.0341 \approx 29$ miljoen jaar.

\textbf{17.} $29 \times \ln 10 \approx 67$ miljoen jaar.

\textbf{18.} $2000\,e^{-0.0341 \times 200} \approx 2000 \times 0.0011 \approx
\qty{2}{m}$.

\textbf{19.} $2000\,e^{-0.0341 \times 250} \approx 2000 \times 0.0002 \approx
\qty{0.4}{m}$: een vlak land. Een top van \qty{2037}{m} betekent dat het model
alleen niet klopt: het gebied moet sindsdien opnieuw zijn opgeheven, of
trager zijn geërodeerd dan het model aanneemt.

\textbf{20.} De gemiddelde hoogte halveert in $29 \times \ln 2 \approx$
\textbf{20 miljoen jaar}.
\end{solution}
